\documentclass[12pt,a4paper]{article}
\usepackage[utf8]{inputenc}
\usepackage[T1]{fontenc}
\usepackage{lmodern}
\usepackage[polish]{babel}
\usepackage{tikz}
\usepackage{geometry}
\geometry{a4paper, margin=2.5cm}

\title{Plan Projektu - MVP}
\author{Tymon Tumialis}
\date{\today}

\begin{document}

\maketitle

\section*{Opis Architektury}

Po przeprowadzeniu analizy (researchu) opracowaliśmy architekturę naszego rozwiązania.

Do statycznej analizy kodu wykorzystamy narzędzie Semgrep z własnymi, zdefiniowanymi antywzorcami oraz przypisanymi do nich kosztami. 

Antywzorce będą przechowywane w pliku YAML. Semgrep przeanalizuje kod użytkownika, a następnie zidentyfikuje i zwróci wszystkie wychwycone antywzorce w pliku JSON. Wynik ten zostanie odczytany przez skrypt w języku Python, którego zadaniem będzie obliczenie sumarycznego kosztu i podjęcie decyzji o ewentualnym uruchomieniu zadania.

W przypadku monitoringu dynamicznego zamierzamy wykorzystać Prolog i Epilog w systemie Slurm, slurmDB, bazę PostgreSQL z rozszerzeniem TimescaleDB oraz środowisko Grafana OS. Schemat architektury przedstawiono poniżej:

\vspace{1cm}
\begin{center}
\begin{tikzpicture}[node distance=2cm, auto]
  \node[draw, rectangle, fill=blue!10, text width=3.5cm, text centered, minimum height=1.5cm, rounded corners] (slurm) {\textbf{SLURM}\\ Prolog i Epilog};
  \node[draw, rectangle, fill=green!10, text width=3.5cm, text centered, minimum height=1.5cm, rounded corners, right of=slurm, node distance=5.5cm] (db) {\textbf{PostgreSQL} \\+ TimescaleDB};
  \node[draw, rectangle, fill=orange!10, text width=3.5cm, text centered, minimum height=1.5cm, rounded corners, right of=db, node distance=5.5cm] (grafana) {\textbf{Grafana OS}};

  \draw[->, thick, >=stealth] (slurm) -- (db);
  \draw[->, thick, >=stealth] (db) -- (grafana);
\end{tikzpicture}
\end{center}
\vspace{1cm}

Monitoring dynamiczny zamierzamy początkowo przetestować w środowisku testowym (sandbox), aby zyskać pewność, że nasze rozwiązanie jest bezpieczne, wydajne i działa poprawnie.

\end{document}
